\documentclass[preprint,12pt,authoryear]{elsarticle}

\usepackage{amsmath,amssymb}
\usepackage{booktabs}
\usepackage{array}
\usepackage{url}
\usepackage{hyperref}
\usepackage{enumitem}
\usepackage{graphicx}
\usepackage{microtype}

\hypersetup{
  colorlinks=true,
  linkcolor=blue,
  citecolor=blue,
  urlcolor=blue
}

\journal{Digital Image Processing — Course Survey}

\begin{document}

\begin{frontmatter}

\title{Deep Learning for Medical Image Segmentation:\\
From Convolutional Networks to Transformers --- A Survey}

\author[a]{Anonymous}
\affiliation[a]{organization={School of Computer Science},
               addressline={Course: Digital Image Processing},
               country={China}}

\begin{abstract}
Medical image segmentation is a cornerstone of modern clinical imaging, enabling
precise delineation of anatomical structures and pathological regions from
modalities such as computed tomography (CT), magnetic resonance imaging (MRI),
and histopathology. Over the past decade, deep learning has fundamentally
transformed this field, supplanting hand-crafted pipelines with end-to-end
trainable architectures that approach or exceed human-level performance on
challenging benchmarks. This survey provides a systematic review of deep
learning methods for medical image segmentation, tracing the evolution from
fully convolutional networks (FCNs) and the landmark U-Net architecture through
attention-augmented convolutional models to the latest Vision Transformer (ViT)
based frameworks. We categorise current approaches into three main families:
(1)~convolutional encoder-decoder networks, covering the U-Net family, DeepLab
variants, and multi-scale feature pyramid methods; (2)~hybrid CNN--Transformer
architectures such as TransUNet and Swin-UNet that inject global context into
local convolutional representations; and (3)~pure 3D transformer models
exemplified by nnFormer~\citep{zhou2023nnformer}, which interleaves convolution
and self-attention with local and global volumetric attention. We also discuss
weakly supervised and semi-supervised learning strategies that alleviate the
high annotation cost characteristic of medical imaging. Key open challenges
including limited labelled data, domain shift, 3D computational cost, and
model explainability are identified, with promising future directions outlined.
This survey is intended as a structured entry point for researchers entering
the field and as a reference for practitioners designing segmentation systems.
\end{abstract}

\begin{keyword}
medical image segmentation \sep deep learning \sep convolutional neural networks
\sep transformer \sep U-Net \sep attention mechanism \sep semantic segmentation
\sep volumetric segmentation
\end{keyword}

\end{frontmatter}

% ─────────────────────────────────────────────────────────────────────────────
\section{Introduction}
\label{sec:intro}

Medical image segmentation---the process of partitioning an image into
semantically meaningful regions---underpins a wide range of clinical workflows,
from tumour volumetry and radiation treatment planning to surgical navigation
and disease progression monitoring~\citep{qureshi2023review}. Unlike natural
scene segmentation, medical images present unique challenges: fine-grained
boundary delineation is required at the voxel level; multiple imaging modalities
(X-ray, CT, MRI, PET, ultrasound) possess heterogeneous noise characteristics;
and obtaining pixel-level ground-truth annotations demands specialist knowledge
that is scarce and expensive.

Classical segmentation pipelines relied on hand-crafted features combined with
statistical or energy-minimisation models. Thresholding, active contours
(``snakes''), graph cuts, and random forests were the dominant paradigms for
two decades. While successful on constrained datasets, these methods generalise
poorly across scanners, acquisition protocols, and pathology types.

The advent of deep learning changed the landscape dramatically. Long~et
al.~\citep{long2015fcn} introduced fully convolutional networks (FCNs),
replacing the fully connected layers of classification networks with
convolutional layers to produce dense pixel-wise predictions. Ronneberger~et
al.~\citep{ronneberger2015unet} subsequently proposed U-Net, whose symmetric
encoder-decoder architecture with skip connections became the \emph{de facto}
standard for biomedical image segmentation.

More recently, transformers---originally designed for natural language
processing~\citep{vaswani2017attention}---have entered computer
vision~\citep{dosovitskiy2020vit} and medical imaging. The self-attention
mechanism's ability to model long-range dependencies across an entire image
or volume addresses a fundamental limitation of convolution, which is
inherently local. Hybrid models such as TransUNet~\citep{chen2021transunet}
embed transformer blocks within U-Net-like encoders, while pure transformer
approaches such as Swin-UNet~\citep{cao2022swinunet} construct full
encoder-decoder pipelines using shifted-window self-attention. For 3D
volumetric data, nnFormer~\citep{zhou2023nnformer} introduces interleaved
convolution and self-attention with both local and global volume-based
attention, significantly outperforming prior transformer baselines on organ
segmentation benchmarks.

\subsection{Paper Organisation}
\label{sec:org}

The remainder of this survey is organised as follows. Section~\ref{sec:protocol}
describes the research protocol. Section~\ref{sec:background} covers background
material including conventional methods and evaluation metrics.
Section~\ref{sec:cnn} reviews CNN-based approaches. Section~\ref{sec:attention}
discusses attention mechanisms within convolutional frameworks.
Section~\ref{sec:transformer} covers transformer-based models.
Section~\ref{sec:weakly} addresses learning under limited annotation.
Section~\ref{sec:discussion} provides discussion including open challenges and
future directions, and Section~\ref{sec:conclusion} concludes.

\subsection{Research Highlights}
\label{sec:highlights}

The main contributions of this review are:
\begin{itemize}[leftmargin=*, itemsep=2pt]
  \item A structured taxonomy of deep learning architectures for medical image
        segmentation covering three generations: FCN/U-Net, attention-augmented
        CNNs, and transformer-based models.
  \item A detailed treatment of nnFormer~\citep{zhou2023nnformer}, the
        state-of-the-art 3D transformer for volumetric segmentation, with
        architectural analysis and benchmark comparison.
  \item A comprehensive discussion of learning under limited annotation,
        including semi-supervised, weakly supervised, and self-supervised
        strategies specific to medical imaging.
  \item Identification of key open challenges (annotation cost, domain shift,
        3D scalability, explainability) with concrete future research directions
        including foundation models and state-space models.
\end{itemize}

% ─────────────────────────────────────────────────────────────────────────────
\section{Research Protocol}
\label{sec:protocol}

\subsection{Research Questions}
\label{sec:rq}

This survey is guided by the following research questions:

\begin{enumerate}[label=RQ\arabic*., leftmargin=*, itemsep=2pt]
  \item How have deep learning architectures for medical image segmentation
        evolved from FCN-based models to transformer-based models?
  \item What are the key architectural innovations that distinguish each
        generation of methods, and how do they address the limitations of
        prior approaches?
  \item Under what conditions do transformer-based models outperform CNN-based
        baselines, and what are the remaining gaps?
  \item What strategies exist for training effective segmentation models when
        labelled data is scarce?
  \item What are the most significant open challenges and promising research
        directions in this field?
\end{enumerate}

\subsection{Comparison with Existing Surveys}
\label{sec:comparison}

Several surveys have reviewed deep learning for medical image segmentation.
Qureshi~et al.~\citep{qureshi2023review} provide a comprehensive review of
semantic segmentation methods for medical imaging, covering FCN-based, region-
based, and weakly supervised methods with extensive application-specific
discussion spanning skin cancer, breast cancer, retinography, brain tumour, and
pulmonary nodules. Our survey complements this work with an in-depth focus on
the transformer revolution post-2020, including systematic analysis of
self-attention mechanisms, hybrid CNN--Transformer architectures, and the latest
3D volumetric transformer nnFormer. Table~\ref{tab:survey_comparison} summarises
the coverage of key topics.

\begin{table}[t]
\caption{Comparison with related survey papers.}
\label{tab:survey_comparison}
\centering
\begin{tabular}{@{}lccccc@{}}
\toprule
\textbf{Survey} & \textbf{CNN} & \textbf{Attention} & \textbf{Transformer}
  & \textbf{3D} & \textbf{Semi-sup.} \\
\midrule
Qureshi et al.~\citeyearpar{qureshi2023review} & \checkmark & \checkmark
  & Partial & -- & \checkmark \\
This survey & \checkmark & \checkmark & \checkmark & \checkmark & \checkmark \\
\bottomrule
\end{tabular}
\end{table}

\subsection{Analysis of Articles}
\label{sec:articles}

The primary sources reviewed in this survey include papers published in leading
venues including IEEE Transactions on Image Processing, IEEE Transactions on
Medical Imaging, Medical Image Analysis, Information Fusion, MICCAI, CVPR,
ICCV, and NeurIPS, covering the period 2015--2024. Search keywords used across
PubMed, IEEE Xplore, and Google Scholar include: \textit{medical image
segmentation deep learning}, \textit{U-Net variants}, \textit{transformer
medical imaging}, \textit{3D segmentation volumetric}, and \textit{semi-
supervised medical segmentation}. The final reference list comprises 21 key
papers selected for their architectural novelty, benchmark impact, and
citation frequency.

% ─────────────────────────────────────────────────────────────────────────────
\section{Background}
\label{sec:background}

\subsection{Conventional Segmentation Methods}
\label{sec:conventional}

Prior to deep learning, medical image segmentation relied on three broad
families of algorithms.

\textbf{Thresholding and region-based methods.} Global or adaptive intensity
thresholding partitions images based on pixel intensity alone. Seeded region
growing (SRG) propagates labels from user-specified seed points to neighbouring
pixels satisfying a homogeneity criterion~\citep{qureshi2023review}. These
methods are computationally efficient but sensitive to imaging noise and
intensity inhomogeneity.

\textbf{Active contour models.} Snakes~\citep{kass1988snakes} represent
boundaries as parametric curves minimising an energy combining internal
smoothness constraints and external image-gradient forces. Level-set methods
reformulate the problem implicitly, enabling topological changes such as contour
splitting. These models provide smooth boundary estimates but require careful
initialisation and hand-tuned parameters.

\textbf{Graph-based methods.} Graph cuts~\citep{boykov2001graphcuts} formulate
segmentation as a minimum-cut problem on a graph where edge weights encode
pairwise boundary costs. These methods produce globally optimal solutions under
their energy models but scale poorly to 3D volumes.

\subsection{Imaging Modalities}
\label{sec:modalities}

Medical images exhibit modality-specific characteristics that influence
algorithm design:
\begin{itemize}[leftmargin=*, itemsep=1pt]
  \item \textbf{CT}: encodes Hounsfield units proportional to X-ray attenuation;
        bone appears bright, soft tissue grey, air dark.
  \item \textbf{MRI}: provides soft-tissue contrast via proton relaxation;
        different pulse sequences (T1, T2, FLAIR) highlight different tissues.
  \item \textbf{Histopathology}: H\&E staining offers cellular-level resolution;
        colour variations encode tissue chemistry.
  \item \textbf{Ultrasound}: suffers from speckle noise and limited contrast;
        real-time acquisition enables interventional guidance.
\end{itemize}

\subsection{Evaluation Metrics}
\label{sec:metrics}

Standard metrics for segmentation evaluation include:

\textbf{Dice Similarity Coefficient (DSC):}
\begin{equation}
  \mathrm{DSC} = \frac{2|P \cap G|}{|P| + |G|}
\end{equation}
where $P$ is the predicted mask and $G$ the ground-truth mask. DSC ranges from
0 (no overlap) to 1 (perfect agreement).

\textbf{Intersection over Union (IoU):}
\begin{equation}
  \mathrm{IoU} = \frac{|P \cap G|}{|P \cup G|}
\end{equation}

\textbf{95th-percentile Hausdorff Distance (HD95):} the 95th percentile of the
directed Hausdorff distance between predicted and ground-truth surface points,
measuring boundary accuracy.

\textbf{Sensitivity and Specificity} quantify true positive and true negative
rates, critical in clinical applications where missed lesions carry different
costs from false positives.

% ─────────────────────────────────────────────────────────────────────────────
\section{CNN-Based Architectures}
\label{sec:cnn}

\subsection{Fully Convolutional Networks}
\label{sec:fcn}

Long~et al.~\citep{long2015fcn} demonstrated that classification networks
(AlexNet, VGG, GoogLeNet) could be converted to dense prediction networks by
replacing fully connected layers with convolutional equivalents and using
bilinear upsampling to restore spatial resolution. The key innovation was skip
connections that fuse semantically rich deep-layer features with spatially
precise shallow-layer features. FCN-8s, using three levels of skip connections,
achieved state-of-the-art results on PASCAL VOC at the time and established the
encoder-decoder paradigm underpinning virtually all subsequent work.

\subsection{U-Net and Its Variants}
\label{sec:unet}

U-Net~\citep{ronneberger2015unet} was specifically designed for biomedical image
segmentation where labelled data is scarce. Its symmetric encoder-decoder uses
max-pooling for downsampling and transposed convolutions for upsampling, with
\emph{concatenation} skip connections (rather than addition) that preserve
high-frequency spatial information in the decoder. Elastic deformation-based
augmentation was proposed to compensate for limited training samples.
Table~\ref{tab:unet_variants} summarises major U-Net variants.

\begin{table}[t]
\caption{Summary of U-Net variants for medical image segmentation.}
\label{tab:unet_variants}
\centering
\begin{tabular}{@{}llp{5cm}@{}}
\toprule
\textbf{Model} & \textbf{Year} & \textbf{Key Contribution} \\
\midrule
U-Net~\citep{ronneberger2015unet} & 2015 &
  Symmetric encoder-decoder with concatenation skip connections \\
3D U-Net~\citep{cciccek2016} & 2016 &
  3D convolutions for volumetric segmentation from sparse annotation \\
U-Net++~\citep{zhou2018unetpp} & 2018 &
  Nested dense skip pathways; deep supervision \\
Attention U-Net~\citep{oktay2018} & 2018 &
  Attention gates on skip connections to focus on relevant regions \\
nnU-Net~\citep{isensee2021nnunet} & 2021 &
  Self-configuring pipeline; SOTA on 23 medical segmentation tasks \\
\bottomrule
\end{tabular}
\end{table}

\textbf{U-Net++}~\citep{zhou2018unetpp} redesigns skip connections as nested,
dense connections at the same resolution, allowing feature aggregation across
multiple semantic scales. Deep supervision is applied at each decoder exit.

\textbf{Attention U-Net}~\citep{oktay2018} places attention gates on skip
connections. Given encoder feature $\mathbf{x}^l$ and decoder gating signal
$\mathbf{g}$:
\begin{equation}
  \alpha^l = \sigma_2\!\left(W_\psi \cdot \sigma_1(W_x \mathbf{x}^l +
  W_g \mathbf{g} + b_g) + b_\psi\right)
\end{equation}
where $\sigma_1$ is ReLU, $\sigma_2$ is sigmoid, and
$\alpha^l \in [0,1]^{H \times W}$ is a spatial attention map applied to skip
features before concatenation. This consistently improves segmentation of small
structures such as the prostate and pancreas.

\textbf{nnU-Net}~\citep{isensee2021nnunet} is a self-configuring framework that
automatically adapts preprocessing, network topology, training, and
post-processing based on dataset fingerprinting. Without manual architecture
search, nnU-Net achieves state-of-the-art performance across a wide range of
tasks and remains a strong competitive baseline against which new methods must
be validated.

\subsection{DeepLab and Multi-Scale Context}
\label{sec:deeplab}

Chen~et al.~\citep{chen2018deeplab} introduced atrous (dilated) convolution,
which enlarges the receptive field without increasing parameter count or
sacrificing spatial resolution. DeepLab v3+ combines an encoder using Atrous
Spatial Pyramid Pooling (ASPP)---parallel atrous convolutions at multiple
rates---with a lightweight decoder that refines object boundaries. This
multi-scale context aggregation is particularly effective for lesions spanning
a wide range of sizes. Feature Pyramid Networks (FPN)~\citep{lin2017fpn}
construct a top-down feature hierarchy with lateral connections, enabling
simultaneous segmentation of structures at all scales.

% ─────────────────────────────────────────────────────────────────────────────
\section{Attention Mechanisms in CNNs}
\label{sec:attention}

\subsection{Channel and Spatial Attention}
\label{sec:channel_spatial}

Squeeze-and-Excitation (SE) networks~\citep{hu2018senet} compute per-channel
attention weights by globally average-pooling spatial dimensions and passing
the result through a bottleneck MLP, allowing the network to recalibrate
feature responses. CBAM (Convolutional Block Attention Module) extends this with
sequential channel-then-spatial attention, improving segmentation in densely
labelled medical volumes.

\subsection{Non-Local Networks}
\label{sec:nonlocal}

Non-Local Networks~\citep{wang2018nonlocal} compute pairwise affinities between
all spatial positions, enabling global context modelling within a convolutional
backbone. The non-local block is inserted as a plug-in module, incurring
$\mathcal{O}(N^2)$ memory cost with spatial resolution $N$. Medical imaging
applications include organ segmentation where global shape constraints matter.

\subsection{Attention Gates in Skip Connections}
\label{sec:attgate}

As described in Section~\ref{sec:unet}, Attention U-Net~\citep{oktay2018}
places spatial attention gates on skip connections. The attention weights are
trained end-to-end via the segmentation loss without additional supervision,
focusing the decoder on task-relevant spatial regions. Ablation studies confirm
consistent 2--5\% DSC improvements on challenging small-organ benchmarks.

% ─────────────────────────────────────────────────────────────────────────────
\section{Transformer-Based Architectures}
\label{sec:transformer}

\subsection{Vision Transformer (ViT)}
\label{sec:vit}

Dosovitskiy~et al.~\citep{dosovitskiy2020vit} demonstrated that a pure
transformer applied to sequences of image patches achieves competitive
classification accuracy when pre-trained on large datasets. An image of size
$H \times W$ is split into $N = HW/P^2$ non-overlapping patches of size
$P \times P$, each linearly projected to a $D$-dimensional embedding. A class
token is prepended, learnable positional encodings are added, and the sequence
is processed by $L$ transformer encoder layers consisting of Multi-head
Self-Attention (MSA) and MLP sublayers. ViT's global receptive field from
the first layer contrasts with the local inductive bias of CNNs.

\subsection{Swin Transformer}
\label{sec:swin}

Liu~et al.~\citep{liu2021swin} introduced the Swin Transformer, addressing
ViT's $\mathcal{O}(N^2)$ attention complexity by computing self-attention
within non-overlapping local windows of size $M \times M$ and shifting the
window partition between successive layers (``shifted window'') to introduce
cross-window connections. This reduces complexity to $\mathcal{O}(N)$ with
respect to image tokens, making Swin practical for high-resolution medical
images. Its hierarchical structure---patch merging reduces spatial resolution
by $2\times$ across four stages---produces multi-scale feature maps compatible
with dense prediction heads.

\subsection{TransUNet: Hybrid CNN-Transformer}
\label{sec:transunet}

Chen~et al.~\citep{chen2021transunet} proposed TransUNet, which uses a ResNet
encoder to extract feature maps, then applies a ViT to the serialised feature
map as a token sequence. The ViT-encoded global context is fed into a U-Net-
like decoder with cascaded upsampling and skip connections. Compared to pure
CNNs, TransUNet achieves improved DSC on multi-organ abdominal CT segmentation,
demonstrating the benefit of combining local convolutional features with global
transformer context.

\subsection{Swin-UNet: Pure Transformer Encoder-Decoder}
\label{sec:swinunet}

Cao~et al.~\citep{cao2022swinunet} built an encoder-decoder entirely from Swin
Transformer blocks, introducing a symmetric Swin decoder with \emph{patch
expanding} operations that upsample feature maps by re-arranging tokens. Skip
connections are retained analogously to U-Net. Swin-UNet achieves better
performance than TransUNet on cardiac MRI and multi-organ CT segmentation,
suggesting that hierarchical Swin features are more suitable than flat ViT
features for dense prediction.

\subsection{nnFormer: 3D Volumetric Transformer}
\label{sec:nnformer}

Zhou~et al.~\citep{zhou2023nnformer} introduced nnFormer (not-another
transFormer), specifically designed for 3D volumetric medical image
segmentation. Three key innovations distinguish nnFormer from prior work.

\textbf{Interleaved Convolution and Self-Attention.} Rather than using
convolutions only in the stem and transformers only at the bottleneck, nnFormer
alternates convolutional blocks with transformer blocks throughout the encoder
and decoder. This interleaving allows each layer to benefit from both local
feature extraction (convolution) and global context (self-attention).

\textbf{Volume-Based Local and Global Attention.} nnFormer introduces two
attention modes: (a)~\emph{local volume-based self-attention} that partitions
the 3D feature map into non-overlapping cubic windows and computes self-
attention within each window; (b)~\emph{global volume-based self-attention}
using dilated windows covering the full volume, capturing inter-region
dependencies. The two modes are applied alternately, analogous to the
local-shifted mechanism in Swin Transformer but generalised to 3D.

\textbf{Skip Attention.} Traditional skip connections concatenate encoder and
decoder features. nnFormer replaces concatenation with a skip attention block
that applies cross-attention between encoder and decoder features, selectively
retaining encoder details most useful for each decoder stage.

Table~\ref{tab:nnformer} summarises nnFormer's performance against
representative baselines on the Synapse multi-organ CT benchmark.

\begin{table}[t]
\caption{Performance on the Synapse multi-organ CT segmentation benchmark
(mean DSC \%).}
\label{tab:nnformer}
\centering
\begin{tabular}{@{}lcc@{}}
\toprule
\textbf{Method} & \textbf{Mean DSC (\%)} & \textbf{HD95 (mm)} \\
\midrule
U-Net~\citep{ronneberger2015unet} & 74.68 & 36.87 \\
TransUNet~\citep{chen2021transunet} & 77.48 & 31.69 \\
Swin-UNet~\citep{cao2022swinunet} & 79.13 & 21.55 \\
nnFormer~\citep{zhou2023nnformer} & \textbf{86.57} & \textbf{10.63} \\
\bottomrule
\end{tabular}
\end{table}

nnFormer achieves an 86.57\% mean DSC, outperforming Swin-UNet by more than
7~percentage points~\citep{zhou2023nnformer}. Its HD95 of 10.63~mm is less than
half that of Swin-UNet (21.55~mm), indicating substantially more accurate
boundary delineation.

% ─────────────────────────────────────────────────────────────────────────────
\section{Learning Under Limited Annotation}
\label{sec:weakly}

\subsection{Semi-Supervised Learning}
\label{sec:ssl}

Semi-supervised methods leverage a small labelled set combined with a large
unlabelled pool. Mean Teacher~\citep{tarvainen2017} maintains an exponential
moving average (EMA) of model weights as a teacher; pseudo-labels from the
teacher guide the student's training on unlabelled data. Uncertainty-guided
filtering selects pseudo-labels with high teacher confidence, avoiding noisy
supervision. Consistency regularisation enforces that strongly augmented views
of the same image produce similar predictions, a principle particularly
valuable in medical imaging where pixel-level annotations are expensive.

\subsection{Weakly Supervised Learning}
\label{sec:wsl}

When only image-level labels or bounding-box annotations are available, Class
Activation Maps (CAMs) from classification networks can generate rough
segmentation seeds~\citep{qureshi2023review}. Scribble-supervised methods accept
sparse curve annotations from radiologists; conditional random field (CRF)
post-processing then propagates labels to unlabelled pixels based on image
gradients, achieving competitive results with a fraction of the annotation cost.

\subsection{Self-Supervised Pre-Training}
\label{sec:selfsup}

Contrastive learning (SimCLR, MoCo) and masked autoencoders (MAE) learn general
visual representations from unlabelled data. Pre-trained models fine-tuned on
small labelled medical datasets consistently outperform random initialisation,
particularly for transformer models requiring large datasets to learn effective
representations. Anatomical priors embedded in pre-training tasks---predicting
relative positions of patches, reconstructing deliberately corrupted
regions---further improve medical imaging performance.

% ─────────────────────────────────────────────────────────────────────────────
\section{Discussion}
\label{sec:discussion}

\subsection{Current Limitations and Challenges}
\label{sec:limitations}

\textbf{Limited and expensive annotation.} Despite progress in semi- and
weakly-supervised learning, most state-of-the-art models still require hundreds
to thousands of annotated cases. The Segment Anything Model
(SAM)~\citep{kirillov2023sam}, trained on over one billion masks from natural
images, offers a promising zero-shot direction but requires domain adaptation to
bridge the gap to medical images. Early adaptations (MedSAM, SAM-Med2D) show
encouraging results.

\textbf{Domain shift and multi-site generalisation.} Intensity distributions
vary across scanners, protocols, and imaging centres (``scanner bias''), causing
models trained at one site to degrade at another~\citep{qureshi2023review}.
Federated learning enables collaborative training without sharing raw patient
data, preserving privacy while improving generalisation.

\textbf{Computational cost of 3D models.} Extending attention to 3D volumes is
memory-intensive. A $128^3$ volume with patch size $4^3$ produces 8,192 tokens;
global self-attention requires $\mathcal{O}(8192^2)$ memory, prohibitive without
special handling. Efficient variants---local windowed attention, linear attention,
low-rank approximation---are active research areas. Hardware-aware design (e.g.,
FlashAttention) reduces memory footprint by avoiding materialisation of full
attention matrices.

\textbf{Explainability and clinical trust.} Deep neural networks remain largely
opaque. Gradient-based attribution (GradCAM, integrated gradients) provides
coarse explanations; attention weights of transformer models offer a more natural
interpretability pathway, though whether attention truly reflects reasoning is
debated. Formal uncertainty quantification (deep ensembles, Monte Carlo dropout)
is critical for flagging unreliable predictions in clinical deployment.

\subsection{Future Directions}
\label{sec:future}

\textbf{Foundation models for medical imaging.} Large vision-language models
accepting text, image, and clinical report inputs simultaneously are emerging.
These models may unify segmentation, detection, and report generation within a
single architecture.

\textbf{Mamba and state-space models.} The Mamba architecture~\citep{gu2023mamba},
based on selective state-space models with linear complexity, has recently been
applied to medical image segmentation (VMamba, SegMamba), offering an alternative
to attention that scales better to high-resolution 3D inputs.

\textbf{Synthetic data and generative augmentation.} Diffusion models and GANs
can synthesise realistic medical images with paired segmentation masks,
alleviating data scarcity in rare pathologies and improving robustness to
imaging variability.

\textbf{Efficient architecture design.} As clinical deployment moves to
resource-constrained environments (edge devices, intraoperative systems), there
is growing need for compact, fast segmentation models that maintain accuracy
without requiring large GPU memory.

% ─────────────────────────────────────────────────────────────────────────────
\section{Conclusion}
\label{sec:conclusion}

This survey has reviewed the evolution of deep learning methods for medical image
segmentation, from the pioneering FCN and U-Net to state-of-the-art transformer-
based architectures. Key milestones include: the encoder-decoder paradigm with
skip connections (U-Net~\citep{ronneberger2015unet}), the self-configuring
pipeline (nnU-Net~\citep{isensee2021nnunet}), hybrid CNN-Transformer designs
(TransUNet~\citep{chen2021transunet}), pure Swin-Transformer encoder-decoders
(Swin-UNet~\citep{cao2022swinunet}), and the volumetric interleaved attention
architecture (nnFormer~\citep{zhou2023nnformer}).

The introduction of transformers has delivered consistent improvements on multi-
organ and cardiac segmentation benchmarks. nnFormer achieves 86.57\% mean DSC
on the Synapse dataset---a 7~percentage-point improvement over Swin-UNet---
confirming the value of combining local convolution and global self-attention
in a unified 3D framework.

Challenges remain in annotation efficiency, domain generalisation, 3D
computational scalability, and clinical explainability. Emerging paradigms
including foundation models, self-supervised pre-training, state-space models,
and diffusion-based augmentation offer compelling paths forward. As the field
converges toward unified architectures that jointly process multi-modal clinical
data, deep learning for medical image segmentation is poised to become a
standard component of clinical imaging workflows.

% ─────────────────────────────────────────────────────────────────────────────
\bibliographystyle{elsarticle-harv}
\bibliography{refs}

\end{document}
